% Chương 3

\chapter{THỰC NGHIỆM} % Tên của chương

\label{Chapter3} % Để trích dẫn chương này ở chỗ nào đó trong bài, hãy sử dụng lệnh \ref{Chapter3} 

%----------------------------------------------------------------------------------------

\section{Kiến trúc tổng quan của hệ thống}

Để giải quyết bài toán nhận diện và đếm phương tiện trong điều kiện khắc nghiệt, đề tài xây dựng và triển khai hệ thống dựa trên một quy trình tuần tự như sau:
\begin{enumerate}
	\item \textbf{Chuẩn bị dữ liệu}: Hệ thống tiếp nhận và phân chia dữ liệu đầu vào thành ba tập riêng biệt: tập Huấn luyện (Training set), tập Xác thực (Validation set) và tập Kiểm thử (Test set).
	
	\item \textbf{Huấn luyện mô hình}: Dữ liệu được nạp vào mạng nơ-ron YOLO11m để trích xuất các đặc trưng của đối tượng.
	
	\item \textbf{Tích hợp vào hệ thống}: Mô hình sau huấn luyện được tích hợp vào luồng xử lý, kết hợp với thuật toán ByteTrack để nhận diện, theo dõi và duy trì mã định danh liên tục cho từng phương tiện.
	
	\item \textbf{Khởi tạo các vùng quan tâm}: Thiết lập các tọa độ của các đa giác trên khung hình video. Các vùng này đóng vai trò là ranh giới để kích hoạt sự kiện đếm.
	
	\item \textbf{Thiết kế logic đếm xe}: Khi một phương tiện nằm trong vùng quan tâm, hệ thống sẽ ghi nhận trạng thái luồng (Đi vào hay đi ra) và tự động tăng biến đếm tương ứng.
	
	\item \textbf{Trực quan hóa}: Hiển thị hộp giới hạn (Bounding box), nhãn phân loại (Class label), mã định danh (Track ID) và thống kê lưu lượng phương tiện đang lưu thông lên video.
	
	\item \textbf{Xuất kết quả}: Video sau quá trình xử lý (đã bao gồm các hiển thị trực quan) được xuất ra và lưu trữ dưới định dạng tiêu chuẩn mp4.
\end{enumerate}

\begin{figure}[htbp]
	\centering
	% height=0.8\textheight: Ép ảnh cao tối đa bằng 80% chiều cao của trang giấy
	% keepaspectratio: Giữ nguyên tỷ lệ, không làm ảnh bị méo hay lùn đi
	\includegraphics[height=0.8\textheight, keepaspectratio]{pipeline.pdf}
	\caption{Sơ đồ kiến trúc tổng quan của hệ thống}
	\label{fig:so_do_tong_quan}
\end{figure}

\section{Huấn luyện mô hình}

\subsection{Chuẩn bị dữ liệu}

Đề tài sử dụng bộ dữ liệu TSBOW (Traffic Surveillance Benchmark for Occluded Vehicles Under Various Weather Conditions) \cite{tsbow} do nhóm nghiên cứu Automation Lab thuộc khoa Kỹ thuật Điện và Máy tính (Department of Electrical and Computer Engineering), Đại học Sungkyunkwan, tỉnh Suwon, Hàn Quốc phát triển và công bố tại hội nghị AAAI-26. Khác với các bộ dữ liệu giao thông thông thường chủ yếu được quay ở góc thấp, TSBOW được thiết kế chuyên biệt cho góc nhìn camera giám sát trên cao. Bộ dữ liệu này tập trung giải quyết những thách thức đến từ sự che khuất phương tiện với mật độ dày đặc và sự suy giảm chất lượng hình ảnh do các hiện tượng thời tiết khắc nghiệt như mưa, tuyết, sương mù và thiếu sáng.

Để phục vụ cho quá trình huấn luyện, toàn bộ dữ liệu huấn luyện và kiểm thử của đề tài được truy xuất và tải về trực tiếp từ kho lưu trữ mã nguồn mở Hugging Face do chính nhóm tác giả cung cấp.

\begin{figure}[H]
	\centering
	
\includegraphics[width=0.8\textwidth]{tsbow.png}
	\caption{Bộ dữ liệu TSBOW.}
	\label{fig:DHTHHN}
\end{figure}

\subsubsection{Thống kê dữ liệu}

TSBOW là một bộ dữ liệu quy mô lớn (Large-scale benchmark), vượt trội về số lượng và chất lượng so với các tập dữ liệu về giao thông \cite{tsbow}. Các thống kê nổi bật bao gồm:
\begin{itemize}
	\item \textbf{Dung lượng video}: Bao gồm 198 video camera giám sát với tổng thời lượng lên tới 32.36 giờ quay thực tế.
	
	\item \textbf{Chất lượng hiển thị}: Toàn bộ dữ liệu được chuẩn hóa ở độ phân giải tiêu chuẩn cao 1280x720 pixel.
	
	\item \textbf{Số lượng khung hình và gán nhãn}: Tập dữ liệu sở hữu tổng cộng 3.2 triệu khung hình. Trong đó, để đảm bảo độ chính xác cho việc đánh giá, có hơn 48.000 khung hình được gán nhãn thủ công (Manually Labeled). Quá trình này đã tạo ra 1.1 triệu hộp giới hạn. Ngoài ra, có tới 71.1 triệu hộp giới hạn được gán nhãn bán tự động để phục vụ cho các mô hình học sâu lớn.
\end{itemize}

\begin{figure}[H]
	\centering
	\includegraphics[width=0.8\textwidth]{statsTable.png}
	\caption{Thống kê dữ liệu TSBOW.}
	\label{fig:DHTHHN}
\end{figure}

\subsubsection{Phân bố nhãn và thuộc tính dữ liệu}

Dữ liệu được thu thập trên 71 tuyến đường và ngã tư khác nhau, mang lại tính đa dạng cao về bối cảnh \cite{tsbow}.

\textbf{Về phân bố nhãn (Classes)}: Các phương tiện được phân loại thành 8 lớp đối tượng: \textbf{Car} \textit{(Ô tô)}, \textbf{Bus} \textit{(Xe buýt)}, \textbf{Truck} \textit{(Xe tải lớn)}, \textbf{Small truck} \textit{(Xe tải nhỏ)}, \textbf{Pedestrian} \textit{(Người đi bộ)}, \textbf{Micromobility} \textit{(Phương tiện nhỏ)}, \textbf{Unidentified} \textit{(Không xác định)}, và \textbf{Others} \textit{(Khác)}. Đặc điểm của phân bố này là sự chênh lệch kích thước vô cùng lớn giữa các phương tiện, đòi hỏi hệ thống phải có khả năng trích xuất đặc trưng đa tỷ lệ \cite{tsbow}.

\textbf{Về thuộc tính dữ liệu (Attributes)}: Môi trường của tập dữ liệu được gán nhãn theo 4 điều kiện thời tiết khác nhau \cite{tsbow}:
\begin{itemize}
	\item \textbf{Thời tiết (Weather)}: Normal (Nắng/Nhiều mây), Haze (Sương mù), Rain (Mưa), và Snow (Tuyết rơi).
	
	\item \textbf{Bối cảnh (Scenario)}: Road (Đoạn đường thẳng), Intersection (Ngã tư giao cắt), Special cases (Trường hợp đặc biệt), và Disaster (Thảm họa thời tiết).
	
	\item \textbf{Đặc điểm đường (Road Type)}: Urban (Đô thị), Standard (Tiêu chuẩn), và Boulevard (Đại lộ).
	
	\item \textbf{Góc nhìn (Scale)}: Fine (Gần/Chi tiết), Medium (Trung bình), và Coarse (Xa/Bao quát)
\end{itemize}

\begin{figure}[H]
	\centering
	\includegraphics[width=0.8\textwidth]{Figure_Pipelines.png}
	\caption{Thuộc tính dữ liệu.}
	\label{fig:DHTHHN}
\end{figure}

\subsection{Môi trường và công cụ thực nghiệm}

\subsubsection{Cấu hình phần cứng}

Quá trình huấn luyện và triển khai được thực hiện trên nền tảng điện toán đám mây Google Colab với bộ xử lý đồ họa (GPU) \textbf{NVIDIA Tesla T4} với 15GB VRAM và bộ xử lý trung tâm (CPU) \textbf{Intel Xeon 2.00GHz} với 12.7GB RAM hệ thống.

\subsubsection{Môi trường phần mềm}

\textbf{Ngôn ngữ lập trình Python 3}: Phiên bản tiêu chuẩn của môi trường Colab.

\textbf{Khung làm việc học sâu (Deep Learning Framework) PyTorch}: Được tích hợp sẵn bộ công cụ NVIDIA CUDA Toolkit để kích hoạt khả năng tăng tốc phần cứng trên GPU.

\textbf{Thư viện triển khai Ultralytics}: Hệ sinh thái mã nguồn mở, được sử dụng để khởi tạo cấu trúc mạng YOLO, tinh chỉnh siêu tham số (Hyperparameters) và tích hợp thuật toán theo dõi quỹ đạo ByteTrack.

\textbf{Thư viện xử lý thị giác OpenCV (Open Source Computer Vision Library)}: Đảm nhiệm việc đọc/ghi luồng video, xử lý từng khung hình, và nội suy đồ họa.

\textbf{Thư viện tính toán số học NumPy}: Được sử dụng để xử lý các phép toán trên ma trận và quản lý mảng dữ liệu.

\subsection{Huấn luyện}

Quá trình huấn luyện đòi hỏi sự cân bằng giữa kích thước đầu vào, độ phức tạp của mô hình và giới hạn phần cứng. Đề tài thiết lập cấu hình huấn luyện tối ưu như sau:
\begin{itemize}
	\item \textbf{Phân chia tập dữ liệu}: Mặc dù bộ dữ liệu TSBOW gốc sở hữu hơn 3.2 triệu khung hình video, tuy nhiên để tối ưu hóa thời gian huấn luyện và tránh hiện tượng mạng YOLO bị học vẹt (Overfitting) do các khung hình liền kề có sự tương đồng với nhau, đề tài sử dụng phiên bản tập dữ liệu đã được trích xuất mẫu (Sampling) với tổng số 25.114 hình ảnh. Toàn bộ dữ liệu đã được nhóm tác giả cấu trúc sẵn thành 3 tập độc lập, với tỷ lệ phân bổ thể hiện trong Bảng~\ref{dataset_split}:
	
%	\begin{table}[h!]
%		\caption{Thông số môi trường thực nghiệm với cấu hình dữ liệu $1000 \times 2048$}
%		\label{tab:env}
%		\centering
%		\begin{tabular}{l l}
%			\toprule
%			\textbf{Thông số} & \textbf{Giá trị} \\
%			\midrule
%			Nền tảng & Google Colaboratory \\
%			GPU & NVIDIA Tesla T4 \\
%			Trình biên dịch & NVCC (CUDA Compiler) \\
%			Ngôn ngữ & CUDA C++ \\
%			Kích thước lát cắt & $1000 \times 2048$ pixels \\
%			Số lượng lát cắt & 500 slices \\
%			Tổng voxel & $\sim 1.024 \times 10^9$ \\
%			\bottomrule
%		\end{tabular}
%	\end{table}
	
	\begin{table}[htbp]
		\centering
		\caption{Cấu hình phân chia tập dữ liệu thực nghiệm}
		\label{dataset_split}
		% Dùng @{} để xóa khoảng trắng thừa ở hai bên mép bảng, giúp bảng khít với lề văn bản
		\begin{tabularx}{\textwidth}{@{} l c c X @{}}
			\toprule
			\textbf{Phân chia dữ liệu} & \textbf{Số lượng} & \textbf{Tỷ lệ} & \textbf{Mục đích sử dụng} \\ 
			\midrule
			Tập huấn luyện & 10.917 & 43,5\% & Trích xuất đặc trưng hình học và cập nhật ma trận trọng số. \\ \addlinespace
			Tập xác thực & 7.583 & 30,2\% & Đánh giá chéo sau mỗi vòng lặp, tính toán sai số để lưu lại mô hình có độ chính xác cao nhất. \\ \addlinespace
			Tập kiểm thử & 6.614 & 26,3\% & Kiểm tra và đánh giá độ chính xác khách quan của mô hình. \\ 
			\bottomrule
		\end{tabularx}
	\end{table}
	
	\item \textbf{Cấu hình siêu tham số (Hyperparameters)}: Quá trình huấn luyện được tinh chỉnh qua các tham số sau:
	\begin{itemize}
		\item[$\circ$] \textbf{Mô hình huấn luyện}: Lựa chọn YOLO11 phiên bản Medium nhằm tối ưu giữa tốc độ và độ chính xác trong quá trình huấn luyện.
		
		\item[$\circ$] \textbf{Kích thước ảnh đầu vào}: Sử dụng ảnh có độ phân giải đầu vào là 1280x1280. Thiết lập này nhằm giữ lại các đặc trưng điểm ảnh của các vật thể nhỏ trong môi trường bị ảnh hưởng bởi thời tiết khắc nghiệt.
		
		\item[$\circ$] \textbf{Chu kỳ huấn luyện}: Quá trình huấn luyện được thiết lập tối đa 50 vòng. Bên cạnh đó còn có cơ chế dừng sớm (Early Stopping), mô hình sẽ dừng huấn luyện nếu chỉ số $mAP$ trên tập Xác thực (Validation) không cải thiện sau 10 vòng liên tiếp. Điều này giúp ngăn chặn hiện tượng quá khớp (Overfitting).
			
		\item[$\circ$] \textbf{Kích thước mẫu huấn luyện}: \texttt{Batch\_size = 8}. Đây là ngưỡng cấu hình an toàn  để khai thác khả năng của GPU mà không gây tràn bộ nhớ.
			
		\item[$\circ$] Ngoài ra còn có một số tham số khác nhằm tối ưu quá trình huấn luyện.
	\end{itemize}
\end{itemize}

\section{Thiết lập Logic theo dõi và đếm xe}

Sau quá trình huấn luyện, đề tài đã trích xuất thành công mô hình tối ưu (best.pt). Mô hình được đưa vào hệ thống để nhận diện, dự đoán hộp giới hạn và phân loại các phương tiện giao thông trong điều kiện thời tiết xấu. Tuy nhiên, để giải quyết bài toán đếm lưu lượng giao thông, hệ thống cần có khả năng theo dõi quỹ đạo chuyển động và thống kê số lượng phương tiện. Toàn bộ quá trình được thiết lập tuần tự thông qua các bước sau:

\subsection{Thiết lập vùng đếm}

Hệ thống trong đề tài sử dụng phương pháp đếm dựa trên vùng quan tâm (Region of Interest - RoI).

Region of Interest (RoI) là một vùng không gian được định nghĩa trong khung hình camera, đại diện cho một khu vực thực tế trên đường. Thay vì chỉ sử dụng một vạch kẻ ngang, việc thiết lập các đa giác (polygon) khép kín cho phép hệ thống thống kê chính xác lưu lượng giao thông theo từng làn đường cụ thể, hỗ trợ luồng giao thông phức tạp.

Để có được tập hợp tọa độ các điểm góc của từng làn đường, đề tài ứng dụng công cụ Roboflow Polygon Annotation. Quy trình thiết lập bao gồm:
\begin{itemize}
	\item Trích xuất một khung hình tĩnh từ video.
	
	\item Tải khung hình lên nền tảng Roboflow và sử dụng công cụ vẽ Polygon để khoanh vùng thủ công các làn đường.
	
	\item Xuất dữ liệu tọa độ $(x, y)$ từ Roboflow và chuyển đổi thành các mảng \texttt{numpy.array} để tích hợp vào mã nguồn hệ thống.
\end{itemize}

\begin{figure}[H]
	\centering
	\includegraphics[width=0.8\textwidth]{polygon.png}
	\caption{Thiết lập vùng đếm.}
	\label{fig:DHTHHN}
\end{figure}

\subsection{Thiết lập logic đếm xe}

Logic đếm xe được thực hiện qua các bước sau:
\begin{itemize}
	\item \textbf{Trích xuất tọa độ và đặc trưng}: Khi một phương tiện di chuyển, hệ thống ghi nhận mã định danh (ID) từ thuật toán ByteTrack, sử dụng điểm trung tâm ở cạnh dưới của hộp giới hạn làm điểm theo dõi $(c_x, c_y)$ và tịnh tiến điểm $c_y$ lên một khoảng bằng 15\% chiều cao hộp giới hạn $h$. Việc này giúp chuẩn hóa tọa độ cho mọi kích thước phương tiện, đảm bảo điểm theo dõi luôn nằm trong vùng đếm. Tham số chiều cao $h$ cũng được lưu lại để phục vụ cơ chế lọc nhiễu.
	
	\item \textbf{Kiểm tra vùng đếm}: Hệ thống lưu trữ một tập hợp (Set) chứa ID của các phương tiện đã được đếm. Nếu ID đã tồn tại trong danh sách này, phương tiện sẽ bị bỏ qua ở các khung hình tiếp theo nhằm tránh lỗi đếm lặp. Ngược lại, với các ID chưa có trong tập hợp, hệ thống sử dụng hàm \texttt{cv2.pointPolygonTest()} để kiểm tra tọa độ $(c_x, c_y)$ có nằm bên trong vùng đếm hay không. Chỉ khi phương tiện tiến vào bên trong vùng đếm, hệ thống mới tiếp tục thực hiện bước tiếp theo.
	
	\item \textbf{Cơ chế lọc nhiễu}: Khi phương tiện tiến vào vùng đếm, hệ thống tính toán khoảng cách không gian (Euclidean distance) giữa phương tiện mới và các phương tiện đã đi qua trong khoảng thời gian 20 khung hình. Hệ thống sẽ thiết lập một ngưỡng sai lệch tối đa bằng 50\% chiều cao $h$ của hộp giới hạn cũ, và được giới hạn trong khoảng từ 30 đến 70 pixel nhằm thích ứng với sự chênh lệch kích thước của phương tiện. Nếu khoảng cách không gian nhỏ hơn ngưỡng giới hạn này, mã định danh mới sẽ được xác định là một ID ảo, từ đó hệ thống bỏ qua việc tăng biến đếm.
	
	\item \textbf{Xác định hướng và ghi nhận đếm}: Khi phương tiện vượt qua được bộ lọc trên, hệ thống đối chiếu tọa độ trục dọc $c_y(t)$ hiện tại với tọa độ cũ $c_y(t-1)$ ở khung hình trước để phân tích hướng di chuyển:
	\begin{itemize}
		\item[$\circ$] \textbf{Nếu $c_y(t) > c_y(t-1)$ (phương tiện di chuyển xuống)}: Tăng biến đếm đi vào.
		
		\item[$\circ$] \textbf{Nếu $c_y(t) < c_y(t-1)$ (phương tiện di chuyển lên)}: Tăng biến đếm đi ra.
	\end{itemize}

\item Sau khi cập nhật, thông số của phương tiện (tọa độ, chiều cao, frame) tiếp tục được sử dụng làm mốc tham chiếu cho bộ lọc, đồng thời ID bị khóa vào tập hợp để chống đếm lặp. Quá trình này được lặp lại cho toàn bộ các phương tiện trên luồng video.
\end{itemize}